% PROGRAMA DE PÓS-GRADUAÇÃO EM ECONOMIA APLICADA
% UNIVERSIDADE FEDERAL DO RIO GRANDE - FURG

% ===========================================================

\documentclass[a4paper, 17pt]{extarticle}

%%%%%%%%%%%%%%%%%%%%%%%%% PACOTES %%%%%%%%%%%%%%%%%%%%%%%%%%%%%%%%%%
 
\usepackage[margin = 0.5in, bottom = 2in]{geometry}
\usepackage[utf8]{inputenc}
\usepackage[T1]{fontenc}
\usepackage{graphicx}
\usepackage{mathrsfs}
\usepackage{enumitem}
\usepackage{multicol}
\usepackage{xparse}
\usepackage{subcaption}
\usepackage{arydshln}
%\usepackage{pdfcomment}
\usepackage{parskip}
\usepackage{xcolor}
\usepackage{lipsum}	
\usepackage{amsfonts,amsmath,amssymb, amsthm, pifont}
\usepackage{booktabs}
\usepackage[colorlinks=true, linkcolor=blue, citecolor=blue, urlcolor=blue]{hyperref} 
\usepackage{todonotes} 
\usepackage{dsfont}
\usepackage{cancel}
\usepackage{layout}
\usepackage[most]{tcolorbox}
\usepackage{tikz}
\usepackage{tikz-3dplot} 
\usepackage[absolute, overlay]{textpos} 
\usetikzlibrary{arrows.meta, cd, patterns, patterns.meta}
\usepackage[bottom]{footmisc}

% Configuração da Fonte
\usepackage{times} % Fonte: Times New Roman

% PACOTES ESSENCIAIS
\usepackage{caption}  
\usepackage{titlesec}  

\titleformat{\section}
  {\normalfont\Large\bfseries} % formato de la fuente
  {\S \hspace{-0cm} \thesection}              % etiqueta de la sección (con símbolo ♦)
  {0.5em}                        % espacio entre el número y el título
  {}                           % código antes del título

%\setlength{\footskip}{10pt} % Espacio entre texto y pie de página
\renewcommand{\footnoterule}{\vspace{5pt}\hrule width 0.3\linewidth\vspace{5pt}}
\setlength{\textheight}{750pt}
%\setlength{\textwidth}{530pt}
%\setlength{\headsep}{10pt}
%\setlength{\marginparwidth}{145pt}
%\setlength{\marginparsep}{12pt}
%\setlength{\topmargin}{-60pt}

\definecolor{verdeoscuro}{RGB}{0,100,0} % Un verde bastante oscuro
\definecolor{rojoscuro}{RGB}{139,0,0} % rojo oscuro tipo "DarkRed"

%\pagestyle{empty}

\newcommand{\ejemplo}[1]{%
  \begin{tcolorbox}[colframe=verdeoscuro, colback=white!1, boxrule=0.8pt]
    #1
  \end{tcolorbox}
}
\newcommand{\alerta}[1]{%
  \begin{tcolorbox}[colframe=red, colback=white!1, boxrule=0.8pt]
    #1
  \end{tcolorbox}
}
\newcommand{\demo}[1]{%
  \begin{tcolorbox}[colframe=gray, colback=white!1, boxrule=0.8pt]
    \textcolor{gray}{#1}
  \end{tcolorbox}
}

\todostyle{orange}{
    linecolor=orange,        % Flecha negra
    bordercolor=orange,      % Borde negro
    textcolor=black!90,      % Borde negro
    backgroundcolor=white!33,  % Sin relleno de color
    size=\scriptsize           % Tamaño de texto
}
\todostyle{green}{
    linecolor=verdeoscuro,        % Flecha negra
    bordercolor=verdeoscuro,      % Borde negro
    textcolor=black!90,      % Borde negro
    backgroundcolor=white!33,  % Sin relleno de color
    size=\scriptsize           % Tamaño de texto
}
\todostyle{red}{
    linecolor=red,        % Flecha negra
    bordercolor=red,      % Borde negro
    textcolor=black!90,      % Borde negro
    backgroundcolor=white!33,  % Sin relleno de color
    size=\scriptsize           % Tamaño de texto
}
\todostyle{red2}{
  linecolor=rojoscuro,        % Flecha negra
  bordercolor=rojoscuro,      % Borde negro
  textcolor=black!90,      % Borde negro
  backgroundcolor=white!33,  % Sin relleno de color
  size=\scriptsize           % Tamaño de texto
}
\todostyle{gris}{
  linecolor=gray,        % Flecha negra
  bordercolor=gray,      % Borde negro
  textcolor=black!90,      % Borde negro
  backgroundcolor=white!33,  % Sin relleno de color
  size=\scriptsize           % Tamaño de texto
}
\begin{document} %Início do Documento

%\layout
%\newpage 


%%%%%%%%%%%%%%%%%%%%%%%%%  MEU PREAMBULO %%%%%%%%%%%%%%%%%%%%%%%%%%%%%%%%%

\makeatletter
\renewcommand\@makefnmark{%
  \hbox{\@textsuperscript{\color{magenta}\normalfont\@thefnmark}}}
\renewcommand\@makefntext[1]{%
  \noindent\makebox[1.8em][r]{\textsuperscript{\color{magenta}\@thefnmark}}#1}
\makeatother

\theoremstyle{definition}
\newtheorem*{definition}{\rotatebox{45}{\large \(\boxdot\)}} 

\theoremstyle{plain}
\newtheorem*{lemma}{{\scriptsize \(\square\)}}
\newtheorem{proposition}{{\large \(\square\)}}
\newtheorem{theorem}{{\large \(\blacksquare\)}} 

\newcommand{\propositionnum}[1]{%
  \renewcommand{\theproposition}{#1}%
}

\newcommand{\theoremnum}[1]{%
\renewcommand{\thetheorem}{#1}%
}
\makeatletter
% Eliminar la puntuación del estilo 'plain' y 'definition'
\def\@thm@headpunct{} % Elimina el punto de los títulos de los teoremas

% Modificar el estilo de 'plain' y 'definition'
\def\th@plain{%
  %\thm@headfont{\normalfont\bfseries} % Mantener la fuente normal, sin negrita
  %\thm@notefont{\normalfont} % El texto de la nota (si la hay) no será en negrita
  \thm@headpunct{} % Sin puntuación adicional 
}

\def\th@definition{%
  %\thm@headfont{\normalfont\bfseries} % Fuente normal para el título
  %\thm@notefont{\normalfont} % El texto de la nota (si la hay) no será en negrita
  \thm@headpunct{} % Sin puntuación adicional
}

\def\th@remark{%
  %\thm@headfont{\normalfont\bfseries} % Fuente normal para el título
  %\thm@notefont{\normalfont} % El texto de la nota (si la hay) no será en negrita
  \thm@headpunct{} % Sin puntuación adicional
}
\makeatother



\theoremstyle{remark}
\newtheorem*{example}{\textcolor{verdeoscuro}{\underline{{Exemplo}}}}
\newtheorem*{exercise}{\textcolor{rojoscuro}{\underline{{Exercise}}}}\theoremstyle{remark}
\newtheorem*{note}{\textnormal{\emph{Nota. }}}
\newtheorem*{corollary}{{\scriptsize \(\boxtimes\)}} 

\newcommand{\id}{\mathds{1}}
\newcommand{\R}{\mathbb{R}}
\newcommand{\F}{\mathbb{F}}
\newcommand{\B}{\mathbb{B}}
\newcommand{\Ei}{}
\newcommand{\Ef}{}
\newcommand{\N}{\mathbb{N}}
\newcommand{\U}{\mathcal{U}}
\newcommand{\V}{\mathcal{V}}
\newcommand{\E}{\mathbb{S}}
\newcommand{\ab}{^{\text{\ding{71}}}}
\newcommand{\ce}{^{\text{\ding{70}}}}
\newcommand{\x}{\textbf{x}}
\newcommand{\y}{\textbf{y}}


\setlength{\columnsep}{1in}

%\begin{multicols}{2}

\-\vspace{-1.3cm} 

\par\noindent \includegraphics[width = \linewidth]{Imagens/Header.png} 

%%%%%%%%%%%%%%%%%%%%%%%% CONFIG. DO TÍTULO %%%%%%%%%%%%%%%%%%%%%%%%

\begin{center} %Inicia a centralização
{%{ \Large QUALI SURVIVAL GUIDE COLLECTION} \\ 
 \LARGE \textbf{\(^{+2}\)Notas de Álgebra Linear II}} \vspace{0.1in}   \\
 { \today }\\ 
 {\small
 \vspace{0.4cm}\emph{"No princípio Deus criou os céus, a terra e...}\hfill Nestor Heli Aponte Avila\(^{1}\) \\ 
\emph{a Álgebra Linear"} - Gênesis 1:1 \hfill \href{mailto:n267452@dac.unicamp.br}{\url{n267452@dac.unicamp.br}}} %\\
%Prof. Piper Pimienta \\
%\href{mailto:pipe@dac.unicamp.br}{\url{pipe@dac.unicamp.br}}}
\end{center} %Finaliza a centralização
\- \vspace{-1cm}

{\small ** Conteúdo basado na disciplina MM719 (Álgebra Linear) ministrada pelo profesor Adriano Moura e  \cite{MA719}, de sua autoría também, no período 2025-II. **\vspace{0.3cm}}%, destinado àqueles que já possuem alguma familiaridade com o assunto, pois não contém demonstrações formais, apenas as ideias por trás delas. Sinceramente, um fã da simplicidade. \\ 


%%%%%%%%%%%%%%%%%%%%%%%%% CONTENIDO %%%%%%%%%%%%%%%%%%%%%%%%%%%%%%%%%%
{\large \textbf{Notação}} \hfill \text{\rotatebox{45}{\large $\boxdot$}}  Definição \hfill \((-)^{\text{\ding{71}}}\)  Aberto \hfill  \((-)^{\text{\ding{70}}}\) Fechado \\ 
{\scriptsize$\square$} \ Lema \hspace{2.4cm} {\large$\square$} \ Proposição \hspace{1.5cm}  {\large $\blacksquare$}  \ Teorema \hspace{2cm} {\scriptsize $\boxtimes$}  \ Corolário
%\begin{center}
%  \setlength{\tabcolsep}{20pt}
%  \begin{tabular}{c c c }
%     \vspace{0.05in}  & &\\ 
%      &   
%  \end{tabular}
%\end{center}

\section*{\S \ \ Preliminares}

É preciso lembrar algumos resultados do curso de Álgebra Linear I, podem ser aprofundados se queser nos capítulos 5 e 6 de \cite{MA719}. Também, algumas propiedades e detalhes do anel de polinômios.  

%\begin{definition}%[]
%    Um par \((A, +)\) é chamado de \emph{grupo abeliano} se \(A\neq \emptyset\) e \(+: A\times A \to A\) é uma operação binaria tal que, \(\forall a, b, c \in A\): 
%    \begin{enumerate}[left = 1.7cm]
%        \item[(G1)] Asocia, \((a+b)+c = a + (a+c)\). 
%        \item[(G2)] Tém neutro, \(\exists ! e \in A : a+e =e+a = a\).  
%        \item[(G3)] Tém inversos, \(\exists ! (-a) \in A : a + (-a)= (-a)+a=e \). 
%        \item[(G4)] Comuta, \(a+b = b + a\). 
%    \end{enumerate}
%\end{definition}

%\begin{definition}
%    Um \emph{corpo} é uma tripla \((\F, +, \cdot)\) onde \(\F\neq \emptyset\) e \(+,\cdot : \F\times \F \to \F\) são operações binarias tais que,
%    \begin{enumerate}[left = 1.7cm]
%        \item[(C1)] \((\F,+)\) é grupo abeliano. 
%        \item[(C2)] O produto \(\cdot\), asocia, comuta e \((\F,\cdot)\setminus \{0\}\) é grupo abeliano.
%        \item[(C3)] Distributividade, \(\forall x,y,z \in \F\), \((x+y)\cdot z= x\cdot z + y\cdot z\). 
%    \end{enumerate}
%\end{definition}

\subsection*{Espaços Vetoriais}

\begin{definition}
   Sejam \(\F\) um corpo, e \(V\neq \emptyset\). Um \(\F\)\emph{-espaço vetorial} é uma tripla \((V,+,\cdot)\) onde, \(+:V\times V \to V\) e \(\cdot : \F \times V \to V\) são operações tais que, \(\forall \lambda, \mu \in \F\), \(\forall v, w \in V\), temos, 
   \begin{enumerate}[left =1.7cm]
        \item[(V1)] \((V,+)\) é grupo abeliano.  
        \item[(V2)] \((V,\cdot)\) asocia, \(\lambda \cdot(\mu \cdot v) = (\lambda \cdot \mu) \cdot v\).  
        \item[(V3)] Distributividade I, \((\lambda + \mu ) \cdot v = \lambda \cdot v + \mu\cdot  v\).  
        \item[(V4)] Distributividade II, \(\lambda \cdot (v+ w) = \lambda \cdot v + \lambda \cdot w\).  
        \item[(V5)] Neutro da multiplicação (por escalar), \(1 \cdot v = v \).  
   \end{enumerate}
\end{definition}

\begin{definition}
    Um subconjunto \(W \subseteq V\) não vacío é dito de \emph{subespaço}, denotamos \(W\leq V\), se \(\forall w_1, w_2 \in W\), \(\forall \lambda \in \F\), temos \( w_1 + \lambda w_2 \in W \).   
\end{definition}

\begin{definition}
    Sejam \(\alpha = (v_i)\) uma familia de vetores em \(V\) e \(m\in \N\). Uma \emph{combinação linear} de \(\alpha\), é um vetor em \(V\) da forma \(v = \sum_{j\leq m} x_{i_j}v_{i_j}\), onde \(x_{i_j} \in \F\).
\end{definition}

\begin{note}
    Denotamos por \([\alpha]\), ao conjunto de todas as combinações lineares de \(\alpha\), repare que \([\alpha]\leq V\). Também, se \(\alpha = \{v\}\) então \([\alpha] = [v] = \F v\). 
\end{note}

\begin{definition}
    Uma familia \(\alpha = (v_i)\) é l.i. sse \(\forall j \in I\), \(v_j \notin [\alpha \setminus \{v_j\}]\). Alem disso, se \([\alpha] = V\) então \(\alpha\) é chamada de \emph{base} de \(V\). 
\end{definition}

\theoremnum{5.5.1.}
\begin{theorem}
    Todo espaço vetorial tém base e quaisquer duas bases tém mesma cardinalidade. \textcolor{gray}{\(\rightarrow \) \cite[Pág. 185]{MA719}}%\textcolor{gray}{\(\rightarrow\) O Teorema da para definir \(\dim V := \#\alpha\)}
\end{theorem}

\begin{definition}
    Seja \(\alpha\) uma base de \(V\). Então, \(\dim V := \# \alpha\) é a \emph{dimensão} de \(V\).  
\end{definition}

\begin{definition}
    Seja \((V_i)\) familia de subespaços em \(V\), definimos a soma deles somo sendo \(\sum V_i := \left[\bigcup V_i\right]\). A soma é direta se \(\forall j\in I\), \(V_j \cap \sum_{i\neq j} V_i =\{0\}\). 
\end{definition}

\begin{note}
    Se a soma é direta escrevemos \(\bigoplus V_i\). 
\end{note}

\propositionnum{5.3.7.}
\begin{proposition}
    \label{prop:5.3.7}
    A soma \(\sum V_i\) é direta \(\leftrightharpoons\) \(\forall v\in \sum V_i\), \(\exists m\in \N \) e \(\exists! v_{i_j} \in V_{i_j}\) tais que \(v = \sum_{j\leq m} v_{i_j}\). \textcolor{gray}{\(\rightarrow\)  \cite[Pág. 174]{MA719}} 
\end{proposition}

\propositionnum{5.4.6.}
\begin{proposition}
    \label{prop:5.4.6}
    Sejam \(\alpha = (v_i)\) e \(V_i = \F v_i\) então \(\alpha\) é l.i. \(\leftrightharpoons\) \(\sum \F v_i\) é direta. 
\end{proposition}

\newcommand{\comentario}[1]{\textcolor{gray}{\(\rightarrow\) #1}}

\begin{corollary}
    \(\alpha\) é base \(\leftrightharpoons\) \(\forall i\in I,\ v_i \neq 0\) e \(V = \bigoplus \F v_i\). Logo, da \(\square \) \ref{prop:5.3.7} temos que \(\forall v\in V\), \(\exists m \in \N\) e \(\exists !x_{i_j} \in \F\) tais que \(v = \sum_{j\leq m } x_{i_j}v_{i_j}\). \comentario{\cite[Pág. 178]{MA719}}
\end{corollary}

\begin{note}
Os \(x_{i_j}\) são as \emph{coordenadas} de \(v\) na base \(\alpha\), identificamos comumente,  
\[(x_{i_j}) =: [v]_\alpha \sim  [ x_{i_1}, \cdots,  x_{i_m}]^T \in M_{m\times 1}(\F).\]
\end{note}
\- \vspace{-1.4cm}
\propositionnum{ \hspace{-0.3cm}}


\subsection*{Transformações Lineares}

\begin{definition}
    Sejam \(V\) e \(W\) \(\F\)-espaços vetoriais. Uma função \(T: V \to W\) é dita \emph{linear} se \(\forall v_1, v_2 \in V\) e \(\forall \lambda \in \F \) tém-se que \(T(v_1 + \lambda v_2) = T(v_1) + \lambda T(v_2)\). 
\end{definition}

\begin{definition}
    Sejam \(T: V \to W\) linear, \(\alpha = (v_j)\) e \(\beta = (w_i)\) bases de \(V\) e \(W \) respetivamente. Então, se \([T(v_j)]_\beta= (a_{ij})\), a \emph{matriz asociada} \([T]_\beta^\alpha := (a_{ij})\) determina \(T\) no sentido que, \(\forall v \in V\), \([T(v)]_\beta = [T]_\beta^\alpha [v]_\alpha\). 
\end{definition}

\begin{note}
    Se \(W= V \) e \(T = \id_{_V}\), então \([\id_{_V}]^\alpha_\beta\) é a matriz cambio de base. 
\end{note}

\theoremnum{6.1.6.}
\begin{theorem}
    \label{thm:6.1.6}
    Sejam \(\alpha = (v_i)\) base de \(V\) e \(\beta = (w_i)\) familia de vetores em \(W\), então \(\exists ! T:V\to W \) linear tal que \(\forall i \in I,\ T(v_i) = w_i\). 
\end{theorem}
    
\begin{note}
    O espaço vetorial das funções \(f: V \to W\), com a soma e o produto usuais é denotado no texto como \(\mathcal{F}(V,W)\).   
\end{note}

\renewcommand{\hom}{\operatorname{Hom}_\F}
\newcommand{\ehom}{\operatorname{End}_\F}

\begin{definition}  
  \(\hom(V,W):= \{T \in \mathcal{F}(V,W): T \text{ é linear}\} \). Se \(V = W\) então denotamos \(\ehom(V) = \hom(V,V)\).\footnote{"\(\operatorname{Hom}\)" vém de homomorfismo e "\(\operatorname{End}\)" de endomorfismo.}   
\end{definition}   

\propositionnum{0.1.}
\begin{proposition}
    %Sejam \(\alpha\) e \(\beta\) bases de \(V\) e \(W\) respetivamente, então, 
\label{prop:0.1}
    Algumas propiedades coletadas de \cite[Seç. 6.1]{MA719}. 
    \begin{enumerate}[left = 0.55cm, label = (\alph*)]
        \item \(\hom(V,W) \leq \mathcal{F}(V,W)\). 
        \item Sejam \(\alpha = (v_j)\) e \(\beta = (w_i) \) bases de \(V\) e \(W\) fixas. A função  \(\zeta : \hom(V,W) \ni T \mapsto [T]^\alpha_\beta \in M_{\#I \times \#J}(\F) \) é linear e injetora. 
        \item Sejam \(\gamma, \ \alpha\) e \(\beta\) bases de \(U,\ V\) e \(W\), \(S \in \hom(U,V)\) e \(T\in \hom(V,W)\) então \(T\circ S \in \hom(U,W)\) e \([T\circ S]^\gamma_\beta = [T]^\alpha_\beta [S]^\gamma_\alpha\).
        \item Se \(\alpha \) e \(\beta \) são bases de \(V\) e \(W\), e \(T \in \hom(V,W)\) é invertivél então \(T^{-1} \in \hom(W,V)\) e \([T^{-1}]^\beta_\alpha = ([T]^\alpha_\beta)^{-1}\).
    \end{enumerate}
\end{proposition}

\begin{definition}
    \(T\in \hom(V,W)\) é um \emph{isomorfismo} se fora sobrejetivo. Dois espaços são \emph{isomorfos}, \(V\cong W\), se \(\exists T\in \hom(V,W)\) isomorfismo.  
\end{definition}

\begin{note}
    A transformação \(\zeta\) do ítem (b) é sobrejetiva se \(\#J < \infty\).
\end{note}

\theoremnum{6.1.9.}
\begin{theorem}
    \(V\cong W\) \(\leftrightharpoons\) \(\dim V = \dim W\). \comentario{\cite[Pág. 191]{MA719}}
\end{theorem}

\propositionnum{\ \hspace{-0.3cm}}
\begin{proposition}
    Seja \(T\in \hom(V,W)\). Se \(U \leq V\) e \(U'\leq W\) então também \(T(U)\leq W\) e \(T^{-1}(U')\leq V\). \comentario{\cite[Pág. 199]{MA719}}
\end{proposition}

\begin{definition}
    \centering 
    \(V_{_T} := T^{-1}(\{0\}) = \mathcal{N}(T)\) \  e \  \(\operatorname{Im}(T) := T(V)\). 
\end{definition}

\begin{note}
    Chamamos estos espaços especiais de \emph{núcleo} e \emph{imagem} de \(T\) e suas dimensões de \emph{nulidade} e \emph{posto}. 
\end{note}

\theoremnum{6.3.6.}
\begin{theorem}
    \(\dim V = \dim V_{_T} + \dim V(T)\). \comentario{\cite[Pág. 201]{MA719}}
\end{theorem}

\propositionnum{6.3.9.}
\begin{proposition}
    \(T \) é injetora \(\leftrightharpoons\) \(V_{_T} = \{0\}\). \comentario{\cite[Pág. 202]{MA719}}
\end{proposition}

\begin{definition}
    Sejam \(\lambda \in \F\) e \(V_\lambda = \{v\in V: T(v) = \lambda v \}\). Se \(\exists v\in V_\lambda \setminus\{0\}\) então \(V_\lambda\) é \emph{autoespaço} e \(v\) um \emph{autovetor}, ambos associados ao \emph{autovalor} \(\lambda\). 
\end{definition}

%\begin{note}
%    No \cite[Pág. 208]{MA719} tém a primeira e equivalente definição de autoespaço como sendo . 
%\end{note}

\begin{definition}
    \(T\) é \emph{diagonalizavél} se \(\exists \beta \) base de \(V\) formada por autovetores.
\end{definition}

\begin{note}
    No caso diagonalizavél e tudo perfeito de mais, pois se \(\lambda_j\) são os autovalores dos \(v_j \in \beta \) então, \([T]^\beta_\beta = \operatorname{diag}(\lambda_j)\). Em versão matricial temos,
    \[[T]^\beta_\beta = [\id_{_V}]^\alpha_\beta [T]^\alpha_\alpha [\id_{_V}]^\beta_\alpha \ \ \sim \ \ B = P^{-1} A P.\]
    É sabido que não todos os operadores são diagonalizavéis. O objetivo do \cite[Cap. 8]{MA719} é ver que, ainda assim, sempre e possivél levar \(T\) a uma matriz quasi-diagonal $B$ (Formas Canônicas). 
\end{note}

\subsection*{O Anel de Polinômios \(\mathcal{P}(\F)\)}

Não vou aprofundar nos detalhes relaçoados com anels, fica muito longe. Por tanto, recomendo rever os conceitos de anel, ideal, domínio de integridade, ideal principal, PID e UFD-Algoritmo de Euclides. 
%\begin{definition}
%    Seja \(R\neq \emptyset \). Um \emph{anel} é uma tripla \((R, +, \cdot) \), com operações tais que, 
%    \begin{enumerate}[left = 1.6cm]
%        \item[(R1)] \((A,+)\) é grupo comutativo.
%        \item[(R2)] \((A, \cdot)\) é asociativa e tém neutro. 
%        \item[(R3)] A multiplicação distribui a soma a esquerda e dereita. 
%    \end{enumerate}
%\end{definition}
\begin{definition}
    Sejam \(I = \{t^k: k\in \mathbb{Z}_\geq0\}\) conjunto respeitando as leis usuais de potências e \((\F,+,\cdot)\) um corpo.\footnote{Quer dizer que \(t^0 =1\) e \(t^m\cdot t^n=t^{m+n}\).  } Defina, \(\forall i \in I\), o espaço vetorial \(V_i = \{a_it^i: a_i \in \F, \  t^i \in I\}\), com as operações, 
    \begin{itemize}
        \item \(+: V_i \times V_i \ni (a_it^i , b_it^i) \mapsto (a_i+b_i)t^i \in V_i\). 
        \item \(\cdot: \F \times V_i \ni (\lambda, a_it^i) \mapsto (\lambda\cdot a_i) t^i \in V_i\). 
%        \item Leis de potências usuais, \(t^0 =1\), \(t^1 = t \) e \(t^k \cdot t^l = t^{k+l}\). 
%        \item Soma, \(\sum^m a_k t^k + \sum^n b_k t^k = \sum^{\max\{n,m\}} (a_k+b_k) t^k\). 
    \end{itemize} 
    Assim, o anel de polinômios com coeficientes em \(\F\) é \(\mathcal{P}(\F) := \sum V_i\).  
\end{definition}

\newcommand{\Pol}{\mathcal{P}(\F)}
\begin{note}
    \(\Pol\) é uma estructura maravilhosa pois, só por mençõar algumas de suas carateristícas que seram de nosso interesse aquí,  
    \begin{itemize}
        \item É espaço vetorial de dimensão infinita e \(\alpha = \{1, t, t^2, \ldots\}\) é uma base, seguindo a \(\square \) \ref{prop:5.3.7}, \(\Pol \ni p(t) = \sum^m_{k=0}a_{i_k}t^{i_k}\). 
        \item Temos um produto bém definido e comutativo que distribui a soma, o que faz dele um anel e mais geralmente uma álgebra comutativa.  
        \item É PID, ou seja que todo ideal é gerado. 
        \item PID \(\Rightarrow \) UFD, ou seja que temos fatorização única em primos, propie-dades de divisibilidade e algoritmo da divisão.
    \end{itemize}
\end{note}

\newcommand{\gr}{\operatorname{gr}}
\begin{definition}
    Seja \(g(t) = \sum^m a_k t^k \in \Pol\) tal que \(a_m \neq 0\). Então o \emph{grau}, \(\gr(g):= m\) e no caso de \(a_m =1\) dizemos que \(g\) é \emph{mônico}. 
\end{definition}
%\begin{definition}
%    Seja \(R\) um anel comutativo e \(I\subseteq R\). Diz-se que \(I\) é um \emph{ideal} se, 
%    \begin{enumerate}[left = 1.4cm]
%        \item[(I1)] \((I,+)\) é um subgrupo aditivo de \((R,+ )\). 
%        \item[(I2)] \(\forall r \in R, \ \forall x \in I\), \(rx \in I\).   
%    \end{enumerate}
%\end{definition}

%\begin{note}
%    \(I\) é \emph{ideal principal} se \(\exists a \in R\) tal que  \(I = \langle a\rangle:= \{ra: r\in R\}\).
%\end{note} 

\newcommand{\dIv}{\operatorname{Div}}
\begin{definition}
    \(\Pol \ni g\) divide \(f \in \Pol\), \(g\ | \ f\),  sse \(\exists q \in \Pol\) tal que \(f = gq\). Também, o conjunto dos divisores é \(\dIv(f):= \{g\in \Pol : g\ | \ f\}\). 
\end{definition}

\begin{definition}
    Diz-se que \(p \in \Pol\setminus \F^\times \) é \emph{primo} sse \(\dIv(p) = \{1, p\}\). 
\end{definition}

\theoremnum{\ \hspace{-0.3cm}}
\begin{theorem}[\emph{Fatoração em primos}]
    Seja \(f\in \Pol \setminus \{1\}\) mônico. Então, \(\{p_j \in \dIv(f): p_j \text{ e primo}\}\) é finito e \(\exists !k_j \in \N\) tais que \(f = p_1^{k_1}p_2^{k_2}\cdots p_n^{k_n}\). 
\end{theorem}

\propositionnum{0.2.}
\begin{proposition}
    \label{prop:0.2}
    Algumas propiedades da divisibilidade, 
    \begin{enumerate}[left = 0.55cm, label = (\alph*)]
        \item \(g \ | \ f\) e \(f \ | \ h \Rightarrow g\ | \ h \).\hspace{3cm} (b) \ \(g \ | \ f \ \leftrightharpoons\  \gr(g) \leq \gr(f)\).   
        \item[(c)] \(g\ | \ f \) e \(g \ | \ h\Rightarrow g\ | \ (f+h)\). \hspace{1.5cm }  (d) \ Se \(g\ | \ f \Rightarrow  \forall r \in \Pol, \   g \ | \ f r\).  \vspace{-0.4cm}\\ 
        \- \hspace{2.7cm}(e) \ \(g \ | \ f\) e \(f \ | \ g \ \leftrightharpoons \ \exists \lambda \in \F^{\times}: f = \lambda g\).  
    \end{enumerate}
\end{proposition}

\theoremnum{ \ \hspace{-0.5cm}}
\begin{theorem}[\emph{Algoritmo da Divisão}]
    Sejam \(f, \ g \in \Pol\) com \(g\neq 0\). Então, \(\exists ! q, r \in \Pol \) tais que \(f = gq +r\) e \( \gr(r) < \gr(g)\). 
\end{theorem}

\newcommand{\mdc}{\operatorname{mdc}}
\theoremnum{\ \hspace{-0.3cm}}
\begin{theorem}
    Sejam \(n\in \mathbb{Z}_{\geq 2}\) e \(f_1, f_2, \ldots, f_n\) não todos nulos, então \(\exists ! d \in \Pol\) mônico ao que chamamos de \(\mdc(f_j)\) que verifica, 
    \begin{enumerate}[label=\roman*., left = 2cm]
        \item \(d \in \bigcap^n \dIv(f_j)\) \hspace{2cm} ii. \ Se \(e \in \bigcap^n \dIv(f_j)  \Rightarrow e \ | \ d\). 
    \end{enumerate}
\end{theorem}

\begin{theorem}[\emph{Bézout}]
    \label{thm:bezout}
    Se \(d = \mdc(f_j)\) então \(\exists a_j \in \Pol\) tais que \(\sum^n a_j f_j = d\). 
\end{theorem}

\begin{note}
    Caso particular \(\mdc(f_1, f_2)=1\), tá dizendo que \(\exists a, b \in \Pol\) tais que \(a(t)f_1(t) + b(t)f_2(t)= 1\), isso sería importante mais pra frente.\footnote{Não falei, mais quando isso acontece diz-se que \(f_1\) e \(f_2\) são coprimos}
\end{note} 
\section{Teoria Geral de Operadores Lineares}

\vspace{-0.3cm}
Se \(T\in \ehom(V)\) então \(T^m := \overbrace{T\circ T \circ \cdots \circ T}^{m\text{ veces }} \in \ehom(V)\), por convenção \(T^0 = \id_{_V}\). Seguindo o ítem (a) da \(\square \) \ref{prop:0.1} , \(\forall p \in \mathcal{P}(\F)\) temos,  
\[p(t) = \sum^m a_kt^k \ \ \ \leadsto \ \ \ P(T) = \sum^m a_kT^k \in \ehom(V).\] 
Também, seguindo a notação dos preliminares \( V_{p}= V_{{p(T)}} = \{v\in V: p(T)(v) = 0\}\). De aquí pra frente seja \(T\in \ehom(V)\) fixo. 

\begin{note}
    Sendo \(p = t-\lambda\), temos uma definição equivalente de autoespaço como sendo \(V_p \neq \{0\}\). 
\end{note}

\subsection{Polinômio Mínimo e Descomposição Primária}

\begin{definition}
    Sejam \(\lambda \in \F\), \(m\in \N\) e \(p(t) = (t-\lambda)^m\). Se \(\exists v\in V_p\setminus \{0\}\) é chamado de \emph{autovetor geralizado} associado ao autovalor \(\lambda\).  
\end{definition}

\begin{definition}
    \(W\leq V\) é \emph{\(T\)-invariante} se \(T(W) \subseteq W\), e a restrição de \(T|_{_W}\) e chamada de o \emph{operador inducido} em \(W\). 
\end{definition}

\begin{lemma}
    \textbf{8.1.1.} Se \(S, T \in \ehom(V)\) são tais que \(S\circ T = T \circ S\), então \(V_{_S}\) é \(T\)-inv. \comentario{ Se \(w \in V_{_S} \Rightarrow S(T(w)) = T(S(w)) = 0 \Rightarrow T(w)\in V_{_S}\)}%ariante.  
\end{lemma}

\begin{note}
    Dado que \(p(T) \circ T = T\circ p(T)\) os subespaçõs \(V_p \) são \(T\)-inv, e também, \(\forall p, f \in \Pol \), \(V_p \subseteq V_{fp}\). Em particular, \(\forall k \in \mathbb{Z}_{\geq 0}, \ V_{p^k} \subseteq V_{p^{k+1}}\). 
\end{note}

\begin{definition}
    Sejam \(\lambda \in F\) e \(V_{\ p}^\infty := \bigcup_{k\geq0} V_{p^k}\). No caso, \(p = (t-\lambda)\) então \(V_{\ p}^\infty \) é chamado de \emph{autoespaço geralizado} associado ao \(\lambda\). 
\end{definition}

\begin{note}
    \(V_{\ p}^\infty \) é \(T\)-inv e se \(\dim V =n < \infty\) então \(V_{\ p}^\infty = V_{p^n}\). 
\end{note}
\begin{definition}
    \centering \(\mathscr{A}_{T}:= \{p \in \Pol: p(T) =0\}\)
\end{definition}

\propositionnum{8.1.3.}
\begin{proposition}
    \label{prop:8.1.3}
    Se \(\dim V = n < \infty \) então \(\mathscr{A}_{T}\neq \{0\}\).  
\end{proposition}

\demo{
Lembre-se que \(\dim (\ehom(V)) = n^2\), logo, \( \exists m \leq n^2\) tal que \(\{T^j\}_{j=0}^m\) é l.d. Tome \(m \) mínimo e faça \(f(T) := T^m - \sum^{m-1} a_k T^k \in \mathscr{A}_{T}\). 
}

\propositionnum{8.1.4.}
\newcommand{\an}{\mathscr{A}}
\begin{proposition}
    Se \(\an_{T}\neq \{0\}\) então \(\exists ! m_T\in \Pol\) mônico tal que \( \forall f \in \an_{T}, \ m_T\ | \ f\).  
\end{proposition}

\demo{  
    \(\an_{T} \) é ideal então \( \exists ! m_T\) tal que \(\an_{T} = \langle m_T\rangle\). Se prefer pegue \(m_T\) de menor grau possivél, suponha \(f = m_T q +r\) e conclua que \(r = 0\). Para unicidade suponha outro e usei \(\square\) \ref{prop:0.2} (e). 
}
\begin{note}
    No caso que \(\an_{T} = \{0\}\) fazemos \(m_T :=0\) é assim que \(V = V_{m_{T}}\). Se \(S = T|_{V_p}\) então \(p \in \an_{_S}\). Também, pode-se verificar que o polinômio construido na \(\square\) \ref{prop:8.1.3} é de fato \(m_T\). 
\end{note}

\hypertarget{lemma816}{}
\begin{lemma}
    \textbf{8.1.6.} Se \(f, g \in \Pol\) co-primos então, a restrição de \(f|_{V_g}\) é injetora.  \comentario{Faz usando o \(\blacksquare\) \hyperref[thm:bezout]{\hspace{0cm}(\emph{Bézout})}, lembre-se que \(1 \leadsto \id_{_V}\)}
\end{lemma}

\propositionnum{8.1.7.}
\begin{proposition}
    \label{prop:8.1.7}
    Seja \(\an_{T} \neq \{0\}\) e \(p \in \Pol\) primo. Então, \(V_p \neq \{0\}\) \(\leftrightharpoons\) \(p\ | \ m_T\). 
\end{proposition}

\demo{
    \((\Rightarrow )\) Suponha \(p \ \nmid \ m_T\) então são co-primos, aplique {\scriptsize\(\square\)} \hyperlink{lemma816}{ 8.1.6.} e conclua \(V_p = \{0\}\). \((\Leftarrow)\) Suponha \(p\ | \ m_T\) e \(V_p = \{0\}\) e contradiga.
} 

\begin{corollary}
    \label{coro:8.1.8}
    \textbf{8.1.8.} Seja \(\an_{T} \neq \{0\}\) e \(p \in \Pol\). Então, \(V_p\neq \{0\}\) \(\leftrightharpoons\) \(\exists f \in \dIv(m_T)\) e \(\exists g \in \Pol\) tais que \(p = gf\). Se \(\nexists q \in \dIv(m_T)\) tal que \(q \ | \ g\) então \(V_p = V_f\). 
\end{corollary}

\demo{
    É coisa de fatorar em primos, ter presentes as propiedades \(\square\) \ref{prop:0.2} (a), \(V_{p}\subseteq V_{fp}\) e conseguir as hipótese da \(\square \) \ref{prop:8.1.7}. 
    %\((\Rightarrow)\) Fatore em primos, \(p = \prod^m p_i^{k_i}\), para \( v \in V_p\setminus \{0\}\), pelo menos um \(p_j^{k_j} \in \an_{_T}\) o anula. Note que \(\{0\} \neq V_{p_j^{k_j}}\supseteq V_{p_j}\) e aplique a \(\square\) \ref{prop:8.1.7}, defina \(f = p_j\). \((\Leftarrow) \) De novo fatore \(f = \prod^m p_i^{k_i}\) pela propiedade \( \square\) \ref{prop:0.2} (a) \(p_j\ | \ m_T\), logo, de novo por \(\square\) \ref{prop:8.1.7} temos \(\{0\}\neq V_{p_j} \subseteq V_f \subseteq V_p\).  
}

\begin{definition}
    Sejam \(v\in V\) e \(\forall k \in \mathbb{Z}_{\geq 0}, \ v_k = T^k(v)\). A familia \(\mathscr{C}^\infty_T(v):= (v_k)\)  é chamada de \(T\)\emph{-ciclo} gerado por \(v\), enquanto \(C_T(v) := [\mathscr{C}^\infty_T(v)]\) é o \emph{subespaço} \(T\)\emph{-cíclico} gerado por \(v\). 
\end{definition}

\begin{note}
    É fácil ver que \(C_T(v)\) é \(T\)-inv. Alternativamente, pode-se definir \(C_T(v) = \mathscr{C}_T^m(v)  \), onde \(m = \dim C_T(v) \in \mathbb{Z}_{\geq 0} \cup \{\infty\}\). 
\end{note}

%\begin{definition}
%    \(W\leq V\) é \(T\)\emph{-cíclico} se \(\exists w \in V\) tal que \(W = C_T(w)\).
%\end{definition}

%\begin{lemma}
%    \(V_\lambda \neq \{0\}\) é \(T\)-cíclico \(\leftrightharpoons\) \(\dim V_\lambda = \{1\}\). 
%\end{lemma}
     
\begin{definition}
    Para \(\alpha = (v_i)\) familia em \(V\) definimos \(C_T(\alpha):= \sum C_T(v_i)\).  
\end{definition}

%\begin{note}
%    Se \(T = \id_{_V}\) ou os \(v_i \in \alpha \) foram autovetores, então \(C_T(\alpha) = [\alpha]\).  
%\end{note}

\begin{note}
    Se \(\dim V = n < \infty\) então num análise analogo da \(\square\) \ref{prop:8.1.3} pode-se construir \(m_{T,v}= m_v\) pensando no mínimo \(m\) tal que \(\mathscr{C}^m_T(v)\) é l.d.. Também, \(C_T(v)\subseteq V_{m_v}\) \(\leftarrow T\)-inv, sendo \(S = T|_{V_{m_v}}\) temos \(m_v = m_S\). 
\end{note}

\begin{corollary}
    \textbf{8.1.9.} \(\forall v \in V\) temos \(m_v \ | \ m_T\). \comentario{Segue da observação acima e \(\square\) \ref{prop:8.1.7}}
\end{corollary}

\propositionnum{8.1.11.}
\begin{proposition}
    \label{prop:8.1.11}
    Sejam \(m \in \N\) e \(p_1, p_2, \ldots, p_m\in \Pol \) co-primos dois a dois. Então a soma \(\sum^m V_{\ p_j}^\infty\) é direta.
\end{proposition}

\demo{
    Seguendo \(\square\) \ref{prop:5.4.6}, pega \(v_j \in V^\infty_{\ p_j}\setminus \{0\}\) e mostra que \(\alpha = (v_j)\) é l.i.. Faz por indução, no paso indutivo suponha \(\sum^{m} a_j v_j = 0 \) e \(g = \prod_{j< m} p_j\). Usa o {\scriptsize\( \square\)} \hyperlink{lemma816}{8.1.6.} com cada \(g_j = p_j\) e \(f = p_m\) pra concluir \(a_m =0\).    
}

\propositionnum{8.1.12.}
\begin{proposition}
    \label{prop:8.1.12}
    Sejam \(m \in \N\) e \(f_1, f_2, \ldots, f_m \in \Pol\) co-primos dois a dois e \(f = \prod^m f_j\). Então, \(V_f = \bigoplus V_{f_j}\). \comentario{Não suporta resumo, olha \cite[Pág. 261]{MA719}}
\end{proposition}

\theoremnum{8.1.13.}
\begin{theorem}[\emph{Descomposição Primária}]
    \label{thm:8.1.13}
    Sejam \(\an_T\neq \{0\}\) e \(\prod^m p_j^{k_j}\) a fatoração em primos de \(m_T\). Então, \(V= \bigoplus V_{p_j^{k_j}}\). 
\end{theorem}

\demo{
   Faça \(p_j^{k_j}= f_j\) na \(\square\) \ref{prop:8.1.12}, o resultado segue do fato de \(V = V_{m_T}\). 
}

\begin{note}
    Os termos da soma são chamados de \emph{subespaços} \(T\)\emph{-primarios}.
\end{note}

\hypertarget{lemma8116}{}
\begin{lemma}
    \textbf{8.1.16.} Sejam \(p\in \Pol\) primo e \(u,v\in V\) tais que \(m_u = m_v = p \). Então, ocurre exatamente uma das duas opções a seguir,
    \begin{enumerate}[label = \roman*., left=2cm]
        \item \(C_T(u) = C_T(v)\). \hspace{2cm} ii. \(C_T(u) \cap C_T(v) = \{0\}\). 
    \end{enumerate} 
\end{lemma}
\begin{itemize}[left=0cm]
    \item[] \(* \ C_T(v) = \{q(T)(v): q \in \Pol\}\ *\) \comentario{Prova meramente conjuntista} 
    \item[] \demo{
    O resultado segue de mostrar \(\forall w \in C_T(v)\), \(C_T(v) = C_T(w)\). Prova as duas contenções usando a igualdade acima e o \(\blacksquare\) \hyperref[thm:bezout]{\hspace{0cm}(\emph{Bézout})}. 
}
\end{itemize}

\propositionnum{8.1.17.}
\begin{proposition}
    Seja \(p\in \Pol\) primo tal que \(\dim V_p <\infty \). Então, \(\exists l\in \mathbb{Z}_{\geq 0}\) e \(v_1, v_2, \ldots, v_l \in V_p\) tais que \(V_p = \bigoplus^l C_T(v_k) \). 
\end{proposition}

\demo{
    Basta ver que \(\forall v\in V_p\setminus \bigoplus^{l-1} C_T(v_k)\) temos \(C_T(v)\cap \bigoplus^{l-1} C_T(v_k) = \{0\}\). Suponha \(w\neq 0 \) naquela interseção e usa os varios fatos em {\scriptsize\(\square\)} \hyperlink{lemma8116}{8.1.16.} pra contradizer, conseguindo ver que \(v \in \bigoplus^{l-1} C_T(v_k)\).  
}

\propositionnum{8.1.18.}
\begin{proposition}
    Seja \(p_j\in \Pol\) primo tal que \(\dim V_{\ p_j}^\infty < \infty\). Então, 
    \[\gr(p_j)\ | \ \dim V_{\ p_j}^\infty \text{  \ \ e \ \  } n_j := \frac{\dim V_{\ p_j}^\infty}{\gr(p_j)} \geq \min\{k: V_{\ p_j}^\infty = V_{p_j^k}\}.\] 
\end{proposition}

\vspace{-0.3cm}
\comentario{
    A demostração não suporta resumo, olha se queser \cite[Pág. 263]{MA719}. 
}

\begin{note}
    Pode e debe-se verificar que \(\min\{k: V_{\ p }^\infty = V_{p^k}\} = k_j\) do \(\blacksquare\) \ref{thm:8.1.13}, é assim que conseguimos definir o \emph{polinômio carateristico} como sendo \(c_T := \prod^m p_j^{n_j} \in \an_T\), exatamente o mesmo do \(\blacksquare\) (\emph{Caley-Hamilton}). 
\end{note}

\ejemplo{
    \centering Encontrando a Descomposição Primária
    \tcblower
    \begin{description}
        \item[Paso 1.] Escolha uma base \(\alpha\) pra seu espaço \(V\) e determine \([T]^\alpha_\alpha =A\). 
        \item[Paso 2.] Encontre \(c_T= \det (t\id_{_V} - A)\) e sua fatoração em primos \(c_T = \prod^m p_j^{n_j}\). \comentario{Precisa habilidade na hora das contas do \(\det\)}
        \item[Paso 3.] Estude subespaços e encontre o \(\min k : \dim V_{p_j^{k}} = \gr(p_j)n_j\).  
    \end{description}
}

\begin{note}
    Na hora das contas são úteis as propiedades,
    \begin{enumerate}[label = \roman*.]
        \item \(\det (a_{ij}) = \sum (-1)^{i+j} a_{ij} M_{ij}\). \comentario{Formula de Laplace}
        \item  {\small\( \det \left(\begin{bmatrix}
            A & 0 \\ 
            C & B
        \end{bmatrix}\right) = \det \left(\begin{bmatrix}
            A & C \\ 
            0 & B
        \end{bmatrix}\right) = \det A \det B\) }. \comentario{\(A\)  e \(B\) quadradas } 
        \item {\small \(\begin{bmatrix}
        | &  & | \\ 
        \mu a_{i1} & \cdots & \mu a_{in}\\
        | &  & |
        \end{bmatrix} = \mu \det \left(\begin{bmatrix}
        | &  & | \\ 
        a_{i1} & \cdots & a_{in}\\
        | &  & |
        \end{bmatrix}\right)\)}. \comentario{Saca múltiplos das filas}
    \end{enumerate}
\end{note}

\subsection{Complementos Invariantes e Bases Cíclicas}

\subsection{Formas Canônicas}

 

%%%%%%%%%%%%%%%%%%%%%%%%%% REFERÊNCIAS %%%%%%%%%%%%%%%%%%%%%%%%%%%%%%%%

\begin{thebibliography}{9}

%\bibitem{martin1966complex}
%Martin, D. y Ahlfors, L.V. (1966). \textit{Complex Analysis}. New York: McGraw-Hill.

\bibitem{MA719}
Moura, Adriano (2025). \textit{Álgebra Linear com Geometría Analítica}. UNICAMP, Campinas, SP. 

\end{thebibliography}

%\end{multicols}

%%%%%%%%%%%%%%%%%%%%%%%%%% APÊNDICES %%%%%%%%%%%%%%%%%%%%%%%%%%%%%%%%

\end{document}
